\chapter*{List of Symbols}

The notation below summarizes the symbols used throughout the mathematical development of the thesis. Where a symbol is overloaded in the source literature, its meaning is determined by the surrounding context.

\begin{longtable}{@{}>{\raggedright\arraybackslash}p{0.22\textwidth} p{0.72\textwidth}@{}}
	\textbf{Symbol}                                & \textbf{Meaning}                                                                                                      \\
	\hline
	\endfirsthead
	\textbf{Symbol}                                & \textbf{Meaning}                                                                                                      \\
	\hline
	\endhead
	$\mathcal{A}$                                  & Auxiliary large language model used for extraction, grouping, and counterfactual generation.                          \\
	$A(C_i,C_j)$                                   & Cramér's $V$ association strength between tabular concepts.                                                           \\
	$\mathcal{M}$                                  & Large language model under evaluation.                                                                                \\
	$M$                                            & Number of concepts in a question.                                                                                     \\
	$\mathbf{X}$                                   & Dataset of questions.                                                                                                 \\
	$X$                                            & Question random variable in the causal graph.                                                                         \\
	$\mathbf{x}$                                   & An observed question from $\mathbf{X}$.                                                                               \\
	$\mathbf{x}_{t\rightarrow t'}$                 & Counterfactual question obtained by replacing the original concept-value tuple $t$ with $t'$.                         \\
	$N$                                            & Number of questions or observations in a dataset; in the generic Shapley formulation, $N$ denotes the player set.     \\
	$N_C$                                          & Total number of concepts across a dataset.                                                                            \\
	$\mathcal{Y}$                                  & Set of possible answer options for a question.                                                                        \\
	$Y$                                            & Categorical random variable representing the answer selected by the evaluated model.                                  \\
	$y$                                            & Realized value of $Y$, i.e., a particular selected answer option with $y\in\mathcal{Y}$.                              \\
	$e$                                            & Natural-language explanation generated by the evaluated model.                                                        \\
	$r=(y,e)$                                      & Model response consisting of an answer and its explanation.                                                           \\
	$\mathbf{C}$                                   & Question-specific set of identified concepts.                                                                         \\
	$C_m$                                          & The $m$th concept in $\mathbf{C}$.                                                                                    \\
	$c_m$                                          & Original value of concept $C_m$.                                                                                      \\
	$c'_m$                                         & Alternative or counterfactual value of concept $C_m$.                                                                 \\
	$\mathbb{C}_m$                                 & Domain of plausible values of concept $C_m$.                                                                          \\
	$\mathbf{K}$                                   & Dataset-wide set of semantic concept categories.                                                                      \\
	$L$                                            & Number of semantic concept categories.                                                                                \\
	$U$                                            & Unobserved state-of-the-world variable in the causal graph.                                                           \\
	$V$                                            & Aspects of a question not represented by its extracted concepts.                                                      \\
	$M_{\Delta}$                                   & Index set of concepts changed by a joint intervention.                                                                \\
	$s$                                            & One correlation group containing concepts that are evaluated jointly.                                                 \\
	$\mathbf{S}$                                   & Partition of a question's concepts into correlation groups.                                                           \\
	$S$                                            & Number of responses sampled per original or counterfactual question.                                                  \\
	$T$                                            & Treatment or coalition, represented as a subset of a correlation group.                                               \\
	$\mathbb{T}$                                   & Cartesian product of the value domains of the concepts in treatment $T$.                                              \\
	$t,t'$                                         & Original and counterfactual joint concept-value tuples.                                                               \\
	$n_{\max}$                                     & Maximum permitted number of concepts in one correlation group.                                                        \\
	$n_{\mathrm{ref}}$                             & Maximum number of treatments included in a reference set.                                                             \\
	$\mathcal{O}(\mathbf{x},\mathbf{C},n_{\max})$  & Oracle that partitions a question's concepts into correlation groups.                                                 \\
	$Z$                                            & Continuous feature vector prior to discretization by the tabular oracle.                                              \\
	$q$                                            & Number of quantile bins used to discretize a continuous feature.                                                      \\
	$Q_r$                                          & Empirical quantile defining a discretization boundary.                                                                \\
	$I_r$                                          & The $r$th interval produced by feature discretization.                                                                \\
	$G=(\mathcal{C},\mathcal{L},W)$                & Weighted association graph used by the tabular oracle, with edge set $\mathcal{L}$.                                   \\
	$v(T)$                                         & Value function measuring the causal effect of treatment $T$.                                                          \\
	$v_{\mathrm{adj}}(T)$                          & Payout-coefficient-adjusted value function.                                                                           \\
	$\phi_m(v)$                                    & Shapley value allocated to concept $C_m$ under value function $v$.                                                    \\
	$\mathrm{CE}(\mathbf{x},C_m)$                  & Scalar Causal Concept Effect for concept $C_m$.                                                                       \\
	$\mathrm{CE}_{\mathrm{raw}}$                   & Scalar unadjusted Shapley Causal Concept Effect.                                                                      \\
	$\mathrm{CE}_{\mathrm{adj}}$                   & Scalar payout-coefficient-adjusted Shapley Causal Concept Effect.                                                     \\
	$\mathrm{CE}_{\mathrm{Matton}}$                & Scalar Causal Concept Effect produced by Matton et al.'s~\cite{mattonWalkTalkMeasuring2024} pipeline.                 \\
	$\mathbf{CE}_{(\cdot)}(\mathbf{x},\mathbf{C})$ & Vector of scalar causal concept effects; the optional subscript distinguishes the raw, adjusted, and Matton variants. \\
	$\mathrm{EE}(\mathbf{x},C_m)$                  & Scalar Explanation-Implied Effect for concept $C_m$.                                                                  \\
	$\mathbf{EE}(\mathbf{x},\mathbf{C})$           & Vector of scalar explanation-implied effects for all concepts in $\mathbf{C}$.                                        \\
	$\mathcal{F}(\mathbf{x})$                      & Faithfulness score for question $\mathbf{x}$.                                                                         \\
	$\mathcal{F}(\mathbf{X})$                      & Dataset-level faithfulness score.                                                                                     \\
	$\mathbb{P}_{\mathcal{M}}$                     & Probability measure induced by the evaluated model $\mathcal{M}$.                                                     \\
	$p_{y\mid I}$                                  & Scalar probability component for answer option $y$ under intervention condition $I$.                                  \\
	$\mathbf{p}_{y\mid I}$                         & Probability vector formed from the scalar components $p_{y\mid I}$ over $y\in\mathcal{Y}$.                            \\
	$D_{\mathrm{KL}}(P\|Q)$                        & Kullback--Leibler divergence from distribution $Q$ to distribution $P$.                                               \\
	$\gamma(\mathbf{x},s)$                         & Payout coefficient of correlation group $s$ for question $\mathbf{x}$.                                                \\
	$\bar{\gamma}$                                 & Reference-set average payout coefficient.                                                                             \\
	$\mathcal{R}$                                  & Reference set used to estimate the average payout coefficient.                                                        \\
	$\kappa(\mathbf{x},n_{\max})$                  & Correlation-strength proxy used in the empirical evaluation.                                                          \\
	$\rho$                                         & Spearman rank-correlation coefficient.                                                                                \\
	$\theta$                                       & Latent model score or logit used in the softmax formulation.                                                          \\
	$\mathcal{E}$                                  & Unobserved stochasticity of the evaluated model.                                                                      \\
\end{longtable}
